\documentclass[11pt]{article}
\usepackage[margin=1in]{geometry}
\usepackage{amsmath,graphicx,booktabs,longtable,hyperref,subcaption}
\graphicspath{{../out/}}

\title{Projecting 3D Human Poses onto 2D Images}
\author{Jeremy}
\date{September 3, 2026}

\begin{document}
\maketitle

\section{The task and the data}
The input is a small Human3.6M-style sample \cite{h36m}: \texttt{data/Pose/poses.txt}
contains 20 non-empty rows, and every row has 45 whitespace-separated values.
The first three values are the camera position $(C_x,C_y,C_z)$ in millimetres.
The remaining 42 values are fourteen $(x,y,z)$ joint positions, also in
millimetres.  The joint order is given by \texttt{joint-names.txt}: Hip,
RHip, RKnee, RAnkle, LHip, LKnee, LAnkle, Neck, LUpperArm, LElbow, LWrist,
RUpperArm, RElbow, and RWrist.  Thus the program treats the input as 20
frames of a 14-joint skeleton, rather than as image-space measurements.

The supplied photographs are 1000 by 1000 pixels.  The challenge also supplies
one scalar focal length, $f=1148.6$, in \texttt{focal.txt}.  For the challenge
projection I use square intrinsics with $f_x=f_y=f$ and principal point
$(c_x,c_y)=(500,500)$ pixels.  The output contains one white-background render,
one photograph overlay, and one side-by-side panel for each of the 20 frames,
plus a generated CSV and LaTeX coordinate table.  Vector and matrix operations
use the header-only GLM library \cite{glm}, while the vendored calibration JSON
is parsed with nlohmann/json \cite{json}.

\section{Identifying the cameras}
The camera position column has only four distinct positions.  Each position
occurs exactly five times, so the 20 frames can be grouped into four camera
views without estimating a new camera for every pose.  Matching is by Euclidean
 distance in millimetres against the centers reconstructed from the vendored
published calibration \cite{calib}, with a tolerance of 1 mm.  The resulting assignment is:

\begin{center}
\begin{tabular}{l l l}
\toprule
Calibration camera & Frames & Center $(x,y,z)$ in mm\\
\midrule
54138969 & 4, 12, 13, 14, 17 & $(1841.107,\ 4955.285,\ 1563.445)$\\
55011271 & 0, 2, 5, 7, 9 & $(1761.279,\ -5078.007,\ 1606.265)$\\
58860488 & 1, 3, 16, 18, 19 & $(-1846.778,\ 5215.047,\ 1491.972)$\\
60457274 & 6, 8, 10, 11, 15 & $(-1794.790,\ -3722.699,\ 1574.893)$\\
\bottomrule
\end{tabular}
\end{center}

For an extrinsic matrix represented as camera-space point
$X_c=RX_w+t$, the center is $C=-R^Tt$.  This is the quantity compared with
the position in each pose row.  The four matches are therefore a direct
identification of the supplied views, not an optimization based on the
skeleton shape.

\section{(a) Choosing a camera orientation}
For the challenge-only orientation, the eye is the row's camera position and
the target is the centroid of its fourteen joints.  With the deterministic
world-up vector $u_w=(0,0,1)$, I construct
\[
 f=\frac{T-E}{\lVert T-E\rVert},\qquad
 r=\frac{f\times u_w}{\lVert f\times u_w\rVert},\qquad
 u=r\times f.
\]
The rotation stores the camera basis as rows: right, negative up, and forward.
The negative-up row makes increasing image $v$ point down, as it does in the
input photographs.  Translation is $t=-RE$, so the complete transform is
$X_c=R X_w-RE$.

A target direction alone determines where the camera looks, but not its roll.
Any rotation of the image plane around $f$ preserves the viewing ray and changes
the projected vertical and horizontal axes.  Choosing the global $z$ axis as
world-up is consequently a convention, not information recovered from the
pose.  It is useful here because it is deterministic and produces a stable
upright image for these views; the implementation rejects the degenerate case
where the view direction is parallel to this axis.  The measured angular
comparison below quantifies the cost of this simple convention relative to the
published orientations.

\section{(b) Projection and the 20 renders}
Every world joint is first transformed to camera coordinates.  Points with
$z\leq0$ are rejected.  For the undistorted pinhole projection, the normalized
coordinates and pixels are
\[
 x_n=\frac{X_c}{Z_c},\qquad y_n=\frac{Y_c}{Z_c},\qquad
 u=f_xx_n+c_x,\qquad v=f_yy_n+c_y.
\]
The look-at mode uses the challenge intrinsics stated above.  The reference
mode uses the matched camera's published $R,t$ and camera-specific intrinsic
matrix.  The command used for these artifacts is \texttt{--mode both}; it
writes the look-at skeleton to the white image and writes the reference
calibration projection over the photograph using raylib's rendering API
\cite{raylib}.  It then combines those two
images into a panel, preserving a direct visual comparison.

\begin{figure}[htbp]
\centering
  \includegraphics[width=0.45\textwidth]{white/00.png}
  \includegraphics[width=0.45\textwidth]{white/01.png}
  \caption{Look-at projected skeletons on a white background (frames 00 and 01).}
\end{figure}

All 20 white renders, overlays, and panels are generated from the same run.
The white images make the coordinate convention and skeleton connectivity
visible without photograph texture; the overlay images test the calibration
projection in the original image coordinate system.

\section{(c) Superimposing on the photographs}
The overlay uses the matched published camera calibration, including its
extrinsic orientation and camera-specific focal/principal-point values.  The
result is composited over the corresponding 1000 by 1000 source frame, and the
panel places that overlay beside the challenge look-at render.  The first panel
is shown below; the remaining 19 panels are generated alongside it in
\texttt{out/panel/}.

\begin{figure}[htbp]
\centering
  \includegraphics[width=0.9\textwidth]{panel/00.png}
  \caption{Frame 00: published-calibration overlay beside the white-background
  look-at projection.}
\end{figure}

\section{Error analysis}
The generated \texttt{out/analysis/joint-error.csv} compares the look-at and
reference pixel coordinates for all 20 frames and all 14 joints.  Its
\texttt{mean\_px} column is the mean of the 14 joint distances for one frame;
\texttt{max\_px} records that frame's largest joint distance.  Averaging the
20 reported frame means gives $97.2925$ pixels, or $97.29$ pixels to the
precision retained in the report.  The worst individual error is $189.88$
pixels, at frame 1, camera 58860488, joint 13 (RWrist).  Frame 1 also has the
largest frame mean, $187.82$ pixels.  These are measured values from the fresh
clean-room output, not values inferred from the images.

The same run writes the spec-required
\texttt{out/analysis/camera-angular-error.csv}.  Averaging its 20 angular
values gives $5.020115^\circ$ ($5.0201^\circ$ rounded); the maximum is
$8.4553^\circ$ at frame 1.  This is the relative rotation angle between the
look-at matrix and the matched published matrix.  The relatively large pixel
error in frame 1 is consistent with its largest orientation discrepancy, but
pixel error also includes differences in intrinsics and in the camera-space
location of each joint.

There is a systematic principal-point and focal-length effect in addition to
orientation.  The look-at construction uses $(500,500)$ and the single
challenge focal, whereas the published cameras have distinct calibration
matrices (their principal points are not all $(500,500)$ and their focal
values are not identical).  A common pixel translation can therefore remain
even when the viewing direction is visually plausible; the perspective scale
also changes with focal length.  The per-camera frame-mean averages illustrate
that the mismatch is view-dependent: 57.09 px for camera 60457274, 78.27 px
for 55011271, 95.95 px for 54138969, and 157.86 px for 58860488.

For this comparison the executable's default is used, so the pinhole
coordinates are not additionally distorted.  The calibration loader retains
the published distortion coefficients and the projection API can apply them,
but enabling distortion would answer a different comparison question and is
not part of the generated default artifacts.

\section{The consolidated coordinate table}
The following table is generated by the executable, rather than copied into
the report.  It contains 20 $\times$ 14 = 280 rows and records frame, matched
camera, joint index/name, and both $(u,v)$ pairs.  Including it with
\verb|\input| keeps the report synchronized with the CSV and makes the exact
rounded coordinates available for inspection.

\input{../out/coords/all-2d-coordinates.tex}

\section{Limitations and references}
The look-at camera is intentionally a simple construction.  It assumes that
the joint centroid is the desired target and that the global $z$ direction is
the correct up convention; neither assumption is learned from the photographs.
The roll is therefore a justified deterministic choice, not a recovered
calibration parameter.  Camera identification also depends on the supplied
positions matching the vendored calibration within the fixed 1 mm tolerance.

Most importantly, \texttt{focal.txt} supplies one focal length for the
challenge projection even though four cameras are present.  Treating 1148.6 as
valid for all four views is required to make the challenge-mode output
reproducible, but it cannot reproduce the camera-specific focal lengths and
principal points in the published calibration.  The reported pixel error is
thus an evaluation of this stated approximation, not a claim that the
look-at model is a replacement for full calibration.  Finally, rendering uses
an off-screen OpenGL context; a headless Linux run needs a provider such as
\texttt{xvfb-run}, while the numerical tests and projection code do not depend
on a display.  The interactive explorer uses raygui for its controls
\cite{raygui}, and the C++ regression suite uses doctest \cite{doctest}.

\begin{thebibliography}{99}
\bibitem{h36m} C.\ Ionescu et al., ``Human3.6M: Large Scale Datasets and
  Predictive Methods for 3D Human Sensing in Natural Environments,''
  IEEE TPAMI, 2014.
\bibitem{calib} K.\ Iskakov, \emph{human36m-camera-parameters}, MIT licence,
  \url{https://github.com/karfly/human36m-camera-parameters}.
\bibitem{raylib} R.\ Santamaria, \emph{raylib}, version 5.5,
  \url{https://www.raylib.com}.
\bibitem{raygui} R.\ Santamaria, \emph{raygui}, version 4.0,
  \url{https://github.com/raysan5/raygui}.
\bibitem{glm} G-Truc Creation, \emph{OpenGL Mathematics (GLM)}, version 1.0.1,
  \url{https://github.com/g-truc/glm}.
\bibitem{json} N.\ Lohmann, \emph{JSON for Modern C++}, version 3.11.3,
  \url{https://github.com/nlohmann/json}.
\bibitem{doctest} V.\ Kirilov, \emph{doctest}, version 2.4.11,
  \url{https://github.com/doctest/doctest}.
\end{thebibliography}

\end{document}
